\documentclass[dvipsnames]{book}

\usepackage[margin=1cm,a5paper]{geometry}

\usepackage{bold-extra}
\usepackage{polski}
\usepackage{iwona}
\usepackage[OT4]{fontenc}

\usepackage{amsmath,amssymb,amsthm}
\usepackage{pgf,tikz}
\usepackage{hyperref}

\pagestyle{empty}
\setlength{\parindent}{0pt}
\setlength{\parskip}{5pt}
\renewcommand\leq\leqslant
\renewcommand\geq\geqslant

\colorlet{maincolor}{ForestGreen}
\usepackage[most]{tcolorbox}
\usepackage{varwidth}

\newtcolorbox{tresc}[3][]{enhanced, sharp corners, boxrule=0mm, colback=maincolor!10,
colframe=black,toptitle=1mm,fonttitle=\bfseries,
colbacktitle=maincolor!25, coltitle=black,
title={Zadanie #2\hfill{\it Opracowanie: #3}},#1}

\newenvironment{problem}[3]{\newpage\begin{tresc}{#1}{#2}{#3}}{\end{tresc}}

% Dodatkowe makra

\usepackage{multicol}
\usepackage{cancel}
\usepackage{bbm}

\def\ind{\mathbbm{1}}

\def\N{\mathbb{N}}
\def\P{\mathbb{P}}
\def\E{\mathbb{E}}
\def\Z{\mathbb{Z}}
\def\Q{\mathbb{Q}}
\def\R{\mathbb{R}}

\def\nwd{\mathrm{NWD}}

\newtheorem*{lemat}{Lemat}

\usepackage{enumitem}
\setlist[enumerate,1]{label={(\alph*)},topsep=0pt}

% ------------------------------------------------------------------
% Kolorowe boxy: definicje (zielone), twierdzenia/obserwacje/wnioski (niebieskie)
% ------------------------------------------------------------------

\colorlet{defcolor}{ForestGreen}
\colorlet{thmcolor}{RoyalBlue}

\newcounter{defcnt}
\newcounter{thmcnt}

% wspólny styl boxa
\tcbset{notatkibox/.style={enhanced, sharp corners, boxrule=0mm, frame hidden,
  toptitle=0.5mm, bottomtitle=0.5mm, left=2mm, right=2mm, top=1.5mm, bottom=1.5mm,
  fonttitle=\bfseries, coltitle=black, breakable, before skip=6pt, after skip=6pt}}

% \begin{definicja}[nazwa]
\NewDocumentEnvironment{definicja}{o}{%
  \refstepcounter{defcnt}%
  \begin{tcolorbox}[notatkibox, colback=defcolor!8, colframe=defcolor!25,
    colbacktitle=defcolor!25,
    title={Definicja \thedefcnt\IfValueT{#1}{\ (#1)}}]%
}{\end{tcolorbox}}

% generator środowisk niebieskich: #1 = nazwa środowiska, #2 = etykieta
\newcommand{\newthmbox}[2]{%
  \NewDocumentEnvironment{#1}{o}{%
    \refstepcounter{thmcnt}%
    \begin{tcolorbox}[notatkibox, colback=thmcolor!8, colframe=thmcolor!25,
      colbacktitle=thmcolor!25,
      title={#2 \thethmcnt\IfValueT{##1}{\ (##1)}}]%
  }{\end{tcolorbox}}%
}

\newthmbox{twierdzenie}{Twierdzenie}
\newthmbox{obserwacja}{Obserwacja}
\newthmbox{wniosek}{Wniosek}
\newthmbox{stwierdzenie}{Stwierdzenie}

% treść zadania -- żółty box, \begin{zadanie}{1.1}
\colorlet{zadcolor}{Dandelion}

\NewDocumentEnvironment{zadanie}{m}{%
  \begin{tcolorbox}[notatkibox, colback=zadcolor!15, colframe=zadcolor!40,
    colbacktitle=zadcolor!40,
    title={Zadanie #1}]%
}{\end{tcolorbox}}

% przykład -- bez boxa, wyróżniony nagłówkiem
\theoremstyle{definition}
\newtheorem*{przyklad}{Przykład}

\renewcommand{\qedsymbol}{$\square$}


\begin{document}

\begin{center}
  {\Large\bfseries Wykład 1 --- 2026-10-02}
\end{center}

\section*{Sprawy organizacyjne}

\begin{itemize}[topsep=0pt,itemsep=0pt]
  \item będą wejściówki,
  \item jedno kolokwium,
  \item egzamin ustny z teorii.
\end{itemize}

\hrulefill

\section*{Proces stochastyczny}

\begin{definicja}[Proces stochastyczny]
$(\Omega, \mathcal{F}, \P)$ --- przestrzeń probabilistyczna, gdzie $\Omega$ jest
niepusty.

$(S, \mathcal{B})$ --- przestrzeń mierzalna, gdzie $S \subseteq \Omega$ oraz
$\mathcal{B} \subseteq \mathcal{F}$ \ (np. $(\R, \mathcal{B}(\R))$, $(\N, 2^{\N})$).

$X \colon (\Omega, \mathcal{F}) \to (S, \mathcal{B})$ --- zmienna losowa.

$T$ --- zbiór indeksów, $T \neq \emptyset$.

Jeżeli $\forall t \in T\ X_t \colon (\Omega, \mathcal{F}) \to (S, \mathcal{B})$,
to $(X_t)_{t \in T}$ nazywamy \emph{procesem stochastycznym} określonym na
$(\Omega, \mathcal{F}, \P)$, o wartościach w $(S, \mathcal{B})$.
\end{definicja}

\begin{przyklad}
Niech $X_1, X_2, X_3, \ldots \overset{\text{iid}}{\sim} \mu$, gdzie $\mu$ jest
rozkładem prawdopodobieństwa na $(\R, \mathcal{B}(\R))$, tzn.
$\forall B \in \mathcal{B}(\R)\ \ \mu(B) \in [0,1]$.

Wtedy $(X_n)_{n \in \N}$ jest procesem stochastycznym --- głupim, bo wszystkie
zmienne losowe są iid, a zatem zachodzi:

$\forall k \in \N\ \forall n_1, \ldots, n_k \in \N\
 \forall A_1, \ldots, A_k \in \mathcal{B}(\R)$:
\[
  \P\big(X_{n_1} \in A_1, \ldots, X_{n_k} \in A_k\big)
  = \prod_{i=1}^{k} \P\big(X_{n_i} \in A_i\big)
  = \prod_{i=1}^{k} \mu(A_i).
\]
\end{przyklad}

\section*{Łańcuchy Markowa z czasem dyskretnym}

\begin{definicja}[Warunek Markowa]
Niech $T = \N_0 := \N \cup \{0\}$, $S$ --- zbiór stanów, co najwyżej
przeliczalny, a $(X_n)_{n \in T}$ --- proces stochastyczny o wartościach w $S$.
Wtedy

$\forall n \in \N\ \ \forall i_0, i_1, \ldots, i_{n-1}, i, j \in S$:
\begin{align*}
  \P\big(X_{n+1} &= j \mid X_0 = i_0, X_1 = i_1, \ldots, X_{n-1} = i_{n-1}, X_n = i\big) \\
  &= \P\big(X_{n+1} = j \mid X_n = i\big),
\end{align*}
gdzie prawa strona to prawdopodobieństwo przejścia w jednym kroku z $i$ do $j$.
Powyższą równość nazywamy \emph{warunkiem Markowa}.
\end{definicja}

\begin{definicja}[Jednorodność łańcucha Markowa]
Łańcuch Markowa jest \emph{jednorodny}, jeśli
\[
  \exists p \colon S^2 \to [0,1] \ \colon\ \ \forall n \in \N\ \ \forall i,j \in S\ \ \
  \P\big(X_{n+1} = j \mid X_n = i\big) = p(i,j).
\]
\end{definicja}

\begin{definicja}[Macierz przejścia i macierz stochastyczna]
$P := \big[p(i,j)\big]_{i,j \in S}$ --- \emph{macierz przejścia}.

$P$ jest macierzą \emph{stochastyczną}, tzn.
\begin{enumerate}[label=\arabic*)]
  \item $\forall i,j\ P_{i,j} \in [0,1]$,
  \item $\forall i \in S\ \sum_{j \in S} P_{i,j} = 1$ \text{(Czyli każdy wiersz
  sumuje się do 1)}.
\end{enumerate}
\end{definicja}

\begin{definicja}[Macierz przejścia w $k$ krokach]
\begin{gather*}
  \forall k \in \N\ \ \exists p_k \colon S^2 \to [0,1]\ \colon\ \
  \forall n \in \N\ \ \forall i,j \in S \\
  \P\big(X_{n+k} = j \mid X_n = i\big) = p_k(i,j).
\end{gather*}
Wtedy $P(k) := \big[p_k(i,j)\big]_{i,j \in S}$ --- \emph{macierz przejścia w $k$
krokach}. Dodatkowo zakładamy, że $P(0) = I$.
\end{definicja}

\begin{stwierdzenie}[Równania Chapmana--Kołmogorowa]
\[
  \forall i,j \in S\ \ \forall m,n \in \N\quad
  p_{m+n}(i,j) = \sum_{k \in S} p_n(i,k) \cdot p_m(k,j).
\]
\end{stwierdzenie}

\begin{proof}
Korzystając z jednorodności:
\begin{align*}
  p_{n+m}(i,j)
  &= \P\big(X_{n+m} = j \mid X_0 = i\big)
   = \frac{\P\big(X_{n+m} = j,\ X_0 = i\big)}{\P\big(X_0 = i\big)} \\[2mm]
  &= \frac{\sum\limits_{k \in S} \P\big(X_{n+m} = j,\ X_n = k,\ X_0 = i\big)}
          {\P\big(X_0 = i\big)} \\[2mm]
  &= \frac{\sum\limits_{k \in S}
      \P\big(X_{n+m} = j \mid X_n = k, \cancel{X_0 = i}\big) \cdot
      \P\big(X_n = k \mid X_0 = i\big) \cdot \cancel{\P\big(X_0 = i\big)}}
          {\cancel{\P\big(X_0 = i\big)}} \\[2mm]
  &= \sum_{k \in S} p_m(k,j) \cdot p_n(i,k),
\end{align*}
gdzie skreślenie warunku $X_0 = i$ wewnątrz $\P$ wynika z $(*)$.
\end{proof}

\begin{wniosek}
\[
  P(n+m) = P(n) \cdot P(m) \implies P(k) = P^k.
\]
\end{wniosek}

\paragraph{$(*)$} Warunek Markowa jest równoważny temu, że
\begin{gather*}
  \forall n \in \N\ \ \forall k \in \{1,2,\ldots,n\} =: [n]\ \
  \forall\, 0 \leq n_1 < n_2 < \ldots < n_k \leq n \\
  \forall x_{n_1}, \ldots, x_{n_k}, x_{n+1} \in S \\
  \P\big(X_{n+1} = x_{n+1} \mid X_{n_1} = x_{n_1}, \ldots, X_{n_k} = x_{n_k}\big)
  = \P\big(X_{n+1} = x_{n+1} \mid X_{n_k} = x_{n_k}\big).
\end{gather*}
Czyli możemy zapomnieć o wszystkich chwilach z przeszłości, z wyjątkiem tej
ostatniej.

\begin{definicja}[Rozkład początkowy]
\[
  \alpha := \big[\P(X_0 = i)\big]_{i \in S}
\]
--- wektor prawdopodobieństwa, \textbf{wierszowy!}
\end{definicja}

\end{document}
